\section[Mengambil Hasil Query (Fetch) dengan MySQLi]{Mengambil Hasil Query (Fetch) dengan MySQLi}

Setelah menjalankan query SELECT, hasilnya berupa object \code{mysqli\_result} yang berisi baris-baris data. Untuk mengakses data, gunakan metode \textit{fetch}: \code{fetch\_assoc()} mengembalikan satu baris sebagai array asosiatif (key = nama kolom), \code{fetch\_row()} sebagai array numerik, \code{fetch\_object()} sebagai object, dan \code{fetch\_all()} mengembalikan semua baris sekaligus. Setelah selesai, bebaskan memori dengan \code{\$result->free()} \cite{phpManual}.

\begin{table}[ht]
\centering
\begin{tabular}{|P{3.5cm}|P{4cm}|P{5cm}|}
\hline
\textbf{Metode Fetch} & \textbf{Return Type} & \textbf{Akses Data} \\
\hline
\code{fetch\_assoc()} & Array asosiatif & \code{\$row['nama']} \\
\hline
\code{fetch\_row()} & Array numerik & \code{\$row[0]} \\
\hline
\code{fetch\_object()} & Object stdClass & \code{\$row->nama} \\
\hline
\code{fetch\_all(MYSQLI\_ASSOC)} & Array of arrays & Semua baris sekaligus \\
\hline
\end{tabular}
\caption{Metode fetch pada MySQLi}
\end{table}

\begin{webcode}[language=PHP, caption={Berbagai metode fetch MySQLi}]
<?php
$result = $conn->query("SELECT id, nama, harga FROM produk");

// fetch_assoc - paling umum
while ($row = $result->fetch_assoc()) {
  echo "Produk: {$row['nama']}, Harga: {$row['harga']}\n";
}

// fetch_all - ambil semua sekaligus (cocok untuk data kecil)
$result = $conn->query("SELECT * FROM kategori");
$semua = $result->fetch_all(MYSQLI_ASSOC);
foreach ($semua as $row) {
  echo "Kategori: {$row['nama']}\n";
}

// Jumlah baris
echo "Total: " . $result->num_rows . " baris\n";
$result->free();
?>
\end{webcode}

